估值结论是本报告最克制的部分：基础经营价值33.08元与现价33.02元几乎相同。最终目标价33.07元由内在价值85\%、市场锚15\%和Broker目标价0\%加权得到，对应约+0.15\%回报。公允价值区间31.64--35.96元用于表达基础倍数不确定性；熊、牛情景11.21/48.34元用于表达完整经营状态，不应混为同一“目标区间”。

\section{完整最终估值表：唯一估值母体为600150.SH}

下表在同一展项内闭合价格、盈利、方法、情景、目标、动作与证据质量。江南、大连、外高桥等经营实体均为合并组成，不是独立估值母体；所有盈利与目标只对应600150.SH合并归母和75.2562亿股。

\begin{exhibitbox}[最终估值总表｜中国船舶（600150.SH）]
\scriptsize
\begin{tabularx}{\exhibitboxwidth}{L{3.2cm}L{4.1cm}L{3.7cm}X}
\toprule
\textbf{估值母体/时点} & \textbf{市场锚} & \textbf{2026--2028 House NP/EPS} & \textbf{主方法} \\
\midrule
唯一parent：600150.SH；A股/CNY；75.2562亿股 & 33.02元@2026-07-22；市值2,484.96亿元 & 2026E 172.76/2.296；2027E 216.49/2.877；2028E 246.81/3.280 & 2027E合并归母周期正常化PE；2.876705元\(\times\)11.5x \\
\bottomrule
\end{tabularx}
\vspace{0.45em}
\begin{tabularx}{\exhibitboxwidth}{L{2.4cm}L{3.8cm}L{4.0cm}X}
\toprule
\textbf{Bear/Base/Bull} & \textbf{最终目标/区间/回报} & \textbf{评级与动作} & \textbf{催化剂} \\
\midrule
11.21/33.08/48.34元 & 33.07元；31.64--35.96元；+0.15\% & 中性偏多（持有/等待验证）；不追价，正式H1后整体重跑 & H1归母/扣非高于101/99亿元，核心GM至少11.72\%，OCF为正，交付接近计划 \\
\bottomrule
\end{tabularx}
\vspace{0.45em}
\begin{tabularx}{\exhibitboxwidth}{L{3.0cm}L{7.4cm}X}
\toprule
\textbf{失效条件} & \textbf{证据质量} & \textbf{零信用边界} \\
\midrule
H1低于92/90亿元；核心GM低于10.5\%；重大延期/撤单/融资/质保；全年OCF/NP低于0.8x且无回收路径 & 中高：R0/R1 maker-checker PASS，价格/股本/盈利/权重可复算；倍数与持续期为House判断；有效Street目标价未披露 & 军品、沪东中华注入、集团未分配订单、弱来源目标价均为0；中远约500亿元项目已在聚合订单内，增量信用0 \\
\bottomrule
\end{tabularx}
\sourcenote{PR-01、CF-03--CF-05、CF-12、CF-14--CF-17、BR-01--BR-13；预测与估值为冻结R1模型。}
\end{exhibitbox}

该总表的“所以呢”是：当前没有由静态估值产生的新增仓位空间。只有催化剂改变合并盈利或现金质量、并使新公允价值相对当时价格形成足够安全边际，动作才升级。

\section{方法选择：2027E周期正常化PE，而非单季年化或船厂SOTP}

2027E是主估值年度，因为2026仍受合并后首次完整年度、交付节奏和H1利润前置影响，2028又对高价订单持续期高度敏感。主方法以2027E Base合并归母EPS 2.876705元乘11.5倍，得到33.0821元。2026E 14倍与2028E 9倍只作期限交叉，OCF/归母用于约束倍数质量。

\begin{exhibitbox}[估值方法与假设桥]
\small
\begin{tabularx}{\exhibitboxwidth}{L{3.0cm}L{4.0cm}L{3.8cm}X}
\toprule
\textbf{项目} & \textbf{选择} & \textbf{理由} & \textbf{失效条件} \\
\midrule
主方法 & 2027E合并归母周期正常化PE & 避免单季/H1年化，匹配周期公司正常化盈利 & 2027盈利分母或持续期显著下修 \\
基础倍数 & 11.5倍 & 位于11.0--12.5倍研究区间中部，接近现价隐含11.48倍 & 毛利/现金质量不达标则去评级 \\
辅助交叉 & 2026E 14倍、2028E 9倍 & 检查期限与周期回落 & 不替代主目标 \\
被阻断方法 & 船型/船厂/军品SOTP、Growth PS/PEG & 逐船经济性与独立EPS不足 & 仅新官方证据和完整模型可解除 \\
现金处理 & 不机械加回1,445.94亿元货币资金 & 未区分受限/经营现金，合同负债反映客户预付款 & 经审计可分配净现金证据 \\
\bottomrule
\end{tabularx}
\sourcenote{CF-03、CF-04；冻结估值模型与审计。}
\end{exhibitbox}

方法选择的“所以呢”是把价值放在一个可复算的合并分母上。只有600150是valuation parent，七家主要船厂不做独立倍数；未披露军品、沪东中华潜在注入、集团未分配订单、邮轮能力期权和弱来源目标价不进入任何情景。

\section{E15｜三情景估值与基础公允区间}

\begin{exhibitbox}[三情景与公允价值]
\small
\begin{tabularx}{\exhibitboxwidth}{L{2.2cm}R{2.2cm}R{2.0cm}R{2.5cm}R{2.4cm}X}
\toprule
\textbf{情景} & \textbf{2027E EPS} & \textbf{PE} & \textbf{价值} & \textbf{相对现价} & \textbf{成立条件} \\
\midrule
熊市 & 1.245290 & 9.0x & 11.21元 & -66.06\% & 交付/毛利断裂并周期去评级 \\
\rowcolor{accentblue!8} 基础 & 2.876705 & 11.5x & 33.08元 & +0.19\% & 按计划兑现、现金逐步改善 \\
牛市 & 3.718618 & 13.0x & 48.34元 & +46.40\% & 交付、结构、毛利和OCF均达上端 \\
\bottomrule
\end{tabularx}
\vspace{0.35em}
\begin{tabularx}{\exhibitboxwidth}{L{4.2cm}R{3.0cm}R{3.0cm}X}
\toprule
\textbf{基础区间计算} & \textbf{低端} & \textbf{高端} & \textbf{用途} \\
\midrule
2027E Base EPS 2.876705元 \(\times\) PE & 11.0x=31.64元 & 12.5x=35.96元 & 表达基础倍数不确定性，不替代完整熊/牛情景 \\
\bottomrule
\end{tabularx}
\sourcenote{冻结R1估值模型；发布数值按两位小数展示。}
\end{exhibitbox}

熊市价值体现盈利与倍数双杀，不能用“在手订单还在”排除；牛市价值要求完整经营上沿，也不能用单项H1利润超预期触发。基础区间覆盖现价，说明当前最合理的行为是等待新证据，而不是将区间上端直接称为目标价。

\section{市场预期桥：现价在支付什么}

按不同观察倍数，33.02元隐含2026/2027/2028归母净利177.50/225.91/276.11亿元。相对House Base，隐含盈利高出2.74\%/4.35\%/11.87\%；所需毛利差随期限扩大，说明市场不仅在支付订单存在，还在支付利润率继续上移和高价订单持续。

\begin{exhibitbox}[市场预期估值桥]
\scriptsize
\begin{tabularx}{\exhibitboxwidth}{L{1.6cm}R{2.0cm}R{1.7cm}X C{1.8cm}R{2.1cm}R{1.9cm}}
\toprule
\textbf{年度} & \textbf{House收入} & \textbf{增速} & \textbf{House归母/EPS} & \textbf{观察倍数} & \textbf{House价值} & \textbf{相对现价} \\
\midrule
2026E & 1,775 & 16.79\% & 172.76/2.296 & 14x & 32.139元 & -2.67\% \\
2027E & 2,000 & 12.68\% & 216.49/2.877 & 11x & 31.644元 & -4.17\% \\
2028E & 2,150 & 7.50\% & 246.81/3.280 & 9x & 29.516元 & -10.61\% \\
\bottomrule
\end{tabularx}
\vspace{0.45em}
\begin{tabularx}{\exhibitboxwidth}{L{2.0cm}R{3.4cm}R{2.2cm}X}
\toprule
\textbf{年度} & \textbf{33.02元隐含归母/EPS} & \textbf{高于House} & \textbf{对应模型要求} \\
\midrule
2026E & 177.50/2.359 & +2.74\% & Growth GM约17.16\%，或归属更优 \\
2027E & 225.91/3.002 & +4.35\% & 综合GM约17.85\%，且持续期不缩短 \\
2028E & 276.11/3.669 & +11.87\% & 综合GM约19.13\%，需要更久高价订单释放 \\
\bottomrule
\end{tabularx}
\sourcenote{现价PR-01；House为冻结增长模型；House价值=House EPS×观察倍数，隐含值为分析推导。金额单位除每股值外为亿元。}
\end{exhibitbox}

现价隐含要求的“所以呢”是：2026小幅超越Base或许足以解释价格，却不足以构成15\%上行；真正的重估需要2027--2028利润持续期显著优于House。若毛利改善停止或现金拖累倍数，表观低远期PE会迅速失去意义。

\section{多锚目标：内在价值85\%、市场15\%、Broker 0\%}

市场锚只使用可观察现价33.02元，不额外资本化军工、集团订单或资产注入叙事。由于没有交易额分位、多周期动量和有效目标价，情绪权重只取15\%；Broker/Street目标价未披露，权重严格为0。

\begin{exhibitbox}[市场隐含情绪锚与最终权重]
\small
\begin{tabularx}{\exhibitboxwidth}{L{3.7cm}R{2.5cm}R{2.0cm}R{3.0cm}X}
\toprule
\textbf{锚} & \textbf{数值} & \textbf{权重} & \textbf{贡献} & \textbf{依据} \\
\midrule
内在/基础价值 & 33.0821元 & 85\% & 28.1198元 & 2027E EPS\(\times\)11.5x \\
市场情绪锚 & 33.02元 & 15\% & 4.9530元 & 可观察现价，有限情绪信用 \\
Broker/Street目标价 & \na & 0\% & 0 & 原始PDF无可用目标价 \\
\midrule
精确多锚值 & \multicolumn{3}{r}{33.072791元} & 发布目标33.07元 \\
\bottomrule
\end{tabularx}
\sourcenote{PR-01、BR-01--BR-13；冻结估值模型。}
\end{exhibitbox}

发布目标涨跌幅按33.07/33.02-1计算为+0.1514\%，约+0.15\%。这个回报与“中性偏多（持有/等待验证）”一致，不能包装为买入。市场锚的作用是避免因0.19\%基础差异机械发出卖出结论，不是为主题热度增加估值。

\section{期限交叉、PEG/PSG与相对估值限制}

2026E按14倍对应32.14元，2028E按9倍对应29.52元，均低于主价值，说明当前价值依赖2027年正常化利润。现价对应House 2026/2027/2028 PE 14.38/11.48/10.07倍，PS 1.400/1.242/1.156倍。PEG/PSG会被低基数和利润率跃升机械压低，只作自身期限观察，不用于主估值。

冻结证据没有业务模式、会计口径和船型结构均匹配的完整同行组，因此不制造“相对韩国/民营船厂折价”结论；PB/ROE因缺少可复算归母权益/BPS而标记未披露。缺失可比并不妨碍绝对估值，却降低相对定价的确定性。

\begin{riskbox}[估值失效与重新定价]
若H1低于92/90亿元、核心毛利低于10.5\%、重大延期/撤单/融资或质保问题出现，或全年OCF/归母低于0.8倍且无回收路径，应切换Bear或降低倍数。若H1高于中点、核心毛利至少11.72\%、OCF为正且交付接近计划，应重跑基本面，而不是简单提高市场锚权重。
\end{riskbox}

估值结论最终回到动作：33.07元目标与现价重合，基础区间没有显著安全边际，Bull又需要多项同步超预期。持有的理由是订单和盈利方向可信；等待的理由是回报不足以补偿持续期与现金风险。
